Each of these differences rests on a single object, a matrix that records which elements of a file produce which of
its features, and on two linear maps through it, one in each direction. Before defining them, we make precise what it
means to explain a detector in one space or the other. Nothing in the construction depends on what an element or a
feature means; we draw our examples from malware detection, where feature extractors are unusually complex.

\subsection{Preliminaries}
\label{sec:XML:Preliminaries}

A classifier cannot read a file directly; the file must first be represented numerically. Following Pierazzi et
al.~\cite{pierazzi2020intriguing}, we call the space of files the \emph{problem space} $\mathcal{Z}$ and the space of
their numerical representations the \emph{feature space} $\mathcal{X}$. A \emph{feature extractor}
$\sigma: \mathcal{Z} \to \mathcal{X}$ moves a file $z$, a sequence of \emph{elements} $z_i$ (here, bytes), to a vector
$x = \sigma(z)$ of \emph{features} $x_j$, and a classifier $\phi: \mathcal{X} \to \mathcal{L}$ assigns that vector a
label. A raw-byte detector such as MalConv is the special case in which $\sigma$ is the identity and the two spaces
coincide. An explainer assigns each input of the function it explains an \emph{attribution}, a real number measuring
that input's contribution to the decision; every explainer in this paper attributes the malware log-odds, a GBDT's raw
score or the difference of a network's two logits.

For a feature-based detector, the inputs an explainer credits depend on where it is applied
(Figure~\ref{fig:feature-vs-problem-explanations}). Applied to the classifier $\phi$, it credits features; applied to
the whole pipeline $\phi \circ \sigma$, extractor included, it credits the elements of the file. The two answers are
explanations of the same decision, but they are different objects, and much of what follows turns on the difference.
The distinction between the two spaces is familiar from problem-space attacks~\cite{pierazzi2020intriguing}, but to our
knowledge it has not been formalized for explanations; we do so here.

\definition[\textit{Feature-Space Explanation}]
\label{def:FeatureSpaceExplanation}
{
For a file $z$ with features $x = \sigma(z)$ and a classifier $\phi$, a feature-space explanation is a vector
$a \in \mathcal{A} = \mathbb{R}^{|x|}$ whose entry $a_j$ is the attribution of feature $x_j$ to $\phi$'s decision.
}

\definition[\textit{Problem-Space Explanation}]
\label{def:ProblemSpaceExplanation}
{
For a file $z$, a feature extractor $\sigma$, and a classifier $\phi$, a problem-space explanation is a vector
$b \in \mathcal{B} = \mathbb{R}^{|z|}$ whose entry $b_i$ is the attribution of element $z_i$ to the decision of
$\phi \circ \sigma$, with $b_i = 0$ for every element that no feature reads.
}

\begin{figure*}[t]
\centering

\tikzset{
    >=Stealth,
    font=\small,
    panelbox/.style={
        rounded corners,
        draw=black!60,
        very thick,
        inner sep=4mm,
        fill=white,
    },
    featurerow/.style={
        draw=black!60,
        rounded corners,
        very thin,
        fill=black!1,
        minimum height=4mm,
        anchor=west
    },
    heatrow/.style={
        draw=black!60,
        rounded corners,
        very thin,
        minimum height=4mm,
        anchor=west
    },
    mininode/.style={
        draw=black,
        thick,
        rounded corners,
        fill=black!2,
        minimum height=0.6cm,
        inner sep=2mm,
        align=center,
        font=\scriptsize
    },
    stage/.style={
        draw=black,
        thick,
        rounded corners,
        minimum width=3cm,
        minimum height=0.9cm,
        align=center,
        fill=black!3
    },
    miniStage/.style={
        draw=black,
        thick,
        rounded corners,
        fill=black!2,
        minimum width=1.0cm,
        minimum height=0.7cm,
        align=center,
        font=\scriptsize
    },
}

\begin{tikzpicture}

% =====================================================
% LEFT PANEL: Feature-space explanation
% =====================================================

% Block containing features + attribution bars
\node[draw=black, thick, rounded corners,
      fill=white,
      minimum width=5.2cm,
      minimum height=2.6cm,
      inner sep=2mm,
      anchor=north west] (fs_block) at (0,0) {};

\begin{scope}[shift={(fs_block.north west)}]
    \node[anchor=south west,font=\scriptsize] at ([xshift=2mm]0,0) {feature $x_j$};
    \node[anchor=south west,font=\scriptsize] at ([xshift=3.4cm]0,0) {attribution $a_j$};
    \node[featurerow, minimum width=3.0cm] (fx1) at ([xshift=2mm,yshift=-4mm]0,0) {\scriptsize \texttt{MaxSectionEntropy}};
    \node[heatrow, minimum width=1.7cm, fill=black!70] (fa1) at ([xshift=3.4cm,yshift=-4mm]0,0) {};
    \node[featurerow, minimum width=3.0cm] (fx2) at ([xshift=2mm,yshift=-10mm]0,0) {\scriptsize \texttt{ImportsHash[17]}};
    \node[heatrow, minimum width=1.2cm, fill=black!45] (fa2) at ([xshift=3.4cm,yshift=-10mm]0,0) {};
    \node[featurerow, minimum width=3.0cm] (fx3) at ([xshift=2mm,yshift=-16mm]0,0) {\scriptsize \texttt{NumStrings}};
    \node[heatrow, minimum width=0.5cm, fill=black!20] (fa3) at ([xshift=3.4cm,yshift=-16mm]0,0) {};
    \node[featurerow, minimum width=3.0cm] (fx4) at ([xshift=2mm,yshift=-22mm]0,0) {\scriptsize \texttt{CheckSum}};
    \node[heatrow, minimum width=1.4cm, fill=black!55] (fa4) at ([xshift=3.4cm,yshift=-22mm]0,0) {};
\end{scope}

% flow diagram (x, phi -> epsilon -> a), directly above the block
\begin{scope}[shift={($(fs_block.north west)+(0.1,2.0)$)}]
    \node[miniStage] (xnode) {$x$};
    \node[miniStage, below=1mm of xnode] (fs_phinode) {$f$};
    \path (xnode) -- (fs_phinode) node[midway] (xphi_mid) {};
    \node[stage, minimum width=2.0cm, align=center, right=5mm of xphi_mid] (fs_epsnode)
    {explainer\\$\epsilon_{f}: \mathcal{X}\mapsto\mathcal{A}$};
    \node[stage, minimum width=2.0cm, align=center, right=3mm of fs_epsnode] (anode)
    {explanation\\$a \in \mathcal{A} \subseteq \mathbb{R}^{|x|}$};
    \draw[->,thick] (xnode.east) -- (fs_epsnode.west);
    \draw[->,thick] (fs_phinode.east) -- (fs_epsnode.west);
    \draw[->,thick] (fs_epsnode.east) -- (anode.west);
\end{scope}

\node[font=\footnotesize, anchor=south, align=center] (fs_title) at ($(fs_block.north)+(0.6,2.85)$)
    {feature-space explanation\\
     $a \in \mathcal{A} \subseteq \mathbb{R}^{|x|}$, here $|x| = 2{,}568$};

\begin{scope}[on background layer]
\node[panelbox, inner sep=2.5mm, fit=(fs_title) (xnode) (anode) (fs_block)] (fs_panel) {};
\end{scope}

% =====================================================
% RIGHT PANEL: Problem-space explanation
% =====================================================

\node[draw=black, thick, rounded corners,
      fill=white,
      minimum width=5.6cm,
      minimum height=2.6cm,
      inner sep=2mm,
      anchor=north west] (ps_block) at ($(fs_block.north east)+(1.1,0)$) {};

\begin{scope}[shift={(ps_block.north west)}]
    \node[anchor=south west,font=\scriptsize] at ([xshift=2mm]0,0) {element(s) $z_i$};
    \node[anchor=south west,font=\scriptsize] at ([xshift=3.9cm]0,0) {attribution $b_i$};
    \node[featurerow, minimum width=3.5cm] (pz1) at ([xshift=2mm,yshift=-4mm]0,0) {\scriptsize \texttt{4D 5A}};
    \node[heatrow, minimum width=0.2cm, fill=black!10] (pb1) at ([xshift=3.9cm,yshift=-4mm]0,0) {};
    \node[featurerow, minimum width=3.5cm] (pz2) at ([xshift=2mm,yshift=-10mm]0,0) {\scriptsize \texttt{.text[0x401000:40]}};
    \node[heatrow, minimum width=1.6cm, fill=black!70] (pb2) at ([xshift=3.9cm,yshift=-10mm]0,0) {};
    \node[featurerow, minimum width=3.5cm] (pz3) at ([xshift=2mm,yshift=-16mm]0,0) {\scriptsize \texttt{.data[0x402000:20]}};
    \node[heatrow, minimum width=0.7cm, fill=black!30] (pb3) at ([xshift=3.9cm,yshift=-16mm]0,0) {};
    \node[featurerow, minimum width=3.5cm] (pz4) at ([xshift=2mm,yshift=-22mm]0,0) {\scriptsize \texttt{overlay}};
    \node[heatrow, minimum width=1.4cm, fill=black!55] (pb4) at ([xshift=3.9cm,yshift=-22mm]0,0) {};
\end{scope}

% flow diagram (z, sigma, phi -> epsilon -> b), directly above the block
\begin{scope}[shift={($(ps_block.north west)+(0.1,2.4)$)}]
    \node[miniStage] (znode) {$z$};
    \node[miniStage, below=1mm of znode] (sigmanode) {$\sigma$};
    \node[miniStage, below=1mm of sigmanode] (phinode) {$f$};
    \node[stage, minimum width=2.0cm, align=center, right=4mm of sigmanode] (epsnode)
    {explainer\\$\epsilon_{\sigma, f}: \mathcal{Z}\mapsto\mathcal{B}$};
    \node[stage, minimum width=2.0cm, align=center, right=3mm of epsnode] (bnode)
    {explanation\\$b \in \mathcal{B} \subseteq \mathbb{R}^{|z|}$};
    \draw[->,thick] (znode.east) -- (epsnode.west);
    \draw[->,thick] (sigmanode.east) -- (epsnode.west);
    \draw[->,thick] (phinode.east) -- (epsnode.west);
    \draw[->,thick] (epsnode.east) -- (bnode.west);
\end{scope}

\node[font=\footnotesize, anchor=south, align=center] (ps_title) at ($(ps_block.north)+(0.4,2.85)$)
    {problem-space explanation\\$b \in \mathcal{B} \subseteq \mathbb{R}^{|z|}$, here $|z| \le 4 \cdot 10^6$};

\begin{scope}[on background layer]
\node[panelbox, inner sep=2.5mm, fit=(ps_title) (znode) (bnode) (ps_block)] (ps_panel) {};
\end{scope}

\end{tikzpicture}

\caption{Two places to explain a feature-based detector. Left: an explainer $\epsilon_{f}$ applied to the classifier attributes the decision to the features. Right: an explainer $\epsilon_{\sigma,f}$ applied to the whole pipeline, extractor included, attributes it to the elements of the file.}
\label{fig:feature-vs-problem-explanations}
\end{figure*}


The two kinds of explanation answer different questions, and an analyst often needs both. A feature-space explanation
says which features a decision rests on, but a feature such as a hashed import bucket or a byte-entropy histogram bin
does not say where in the file it comes from. A problem-space explanation says which bytes matter, but not what the
detector computed from them. For a raw-byte detector, any explainer gives a problem-space explanation directly and no
feature-space one at all. For a feature-based detector, the reverse holds: a feature-space explanation is routine,
while a problem-space explanation is hard to obtain. The extractor is neither differentiable nor invertible,
so gradient-based explainers do not apply; a perturbation-based explainer must keep every perturbed file parseable, and
must re-run the extractor, which parses the file, for each of its queries. Both difficulties disappear if an
explanation computed in the space where it is easy can be moved to the space where it is needed.

\subsection{The Provenance Matrix}
\label{sec:XML:ProvenanceMatrix}

Some features are produced by no element of the file. A few are computed from no particular element, such as a flag
recording whether the file could be parsed; many more take their value because something is absent, such as a hashed
import bucket into which no imported function falls. In hashed and structural feature sets, such features can be the
majority. To account for them without fabricating influence, we augment every file with a \emph{null element}
$z_\varnothing$, which stands for ``no part of $z$'', and write $\bar{z} = (z_1, \dots, z_{|z|}, z_\varnothing)$.

\definition[\textit{Explanation Provenance}]
\label{def:ExplanationProvenance}
{
The explanation provenance of a file $z$ with features $x = \sigma(z)$ is a matrix
$P \in \mathbb{R}^{(|z| + 1) \times |x|}$ whose entry $P_{i,j}$ is the share of feature $x_j$ that element $\bar{z}_i$
produces, subject to four axioms:
\begin{itemize}
    \item \textit{Nonnegativity:} $P_{i,j} \ge 0$ for all $i, j$.
    \item \textit{Normalization:} $\sum_i P_{i,j} = 1$ for all $j$, the sum including the null element.
    \item \textit{Influence:} for every element $z_i$ of $z$, $P_{i,j} > 0$ if and only if $z_i$ influences $x_j$.
    \item \textit{Null:} $P_{\varnothing,j} = 1$ if no element of $z$ influences $x_j$, and $0$ otherwise.
\end{itemize}
}

The definition leaves open what it means for an element to influence a feature, and how much of the feature it
produces. Both are fixed when $P$ is constructed for a particular extractor, and Section~\ref{sec:XML:Constructing}
states the choices we make for EMBER.

We call the features that no element influences \emph{unattributable} and write $U = \{ j : P_{\varnothing,j} = 1 \}$.
Without the null element, an unattributable feature could not satisfy both normalization and influence. Dropping it
from $P$ would break normalization, and the transfers would lose its attribution; spreading it uniformly over the
elements would fabricate influence and smear its attribution over the whole file. The null element keeps every axiom,
is simply one more row of $P$, and makes the attribution that no part of the file can carry an explicit quantity.

\subsection{Pullback and Pushforward}
\label{sec:XML:Transfers}

The provenance matrix transfers explanations between the two spaces, in both directions
(Figure~\ref{fig:pullback-pushforward}).

\definition[\textit{Pullback}]
\label{def:Pullback}
{
The pullback of a feature-space explanation $a \in \mathcal{A}$ is the problem-space explanation
$\tau(a) = Pa = b \in \mathbb{R}^{|z| + 1}$ of $\bar{z}$. Its last entry, $b_\varnothing = \sum_{j \in U} a_j$, is the
attribution that no element of $z$ carries.
}

Each feature's attribution is split among the elements that produce it, in proportion to their shares. Because every
column of $P$ sums to one, the pullback conserves attribution, $\sum_i b_i = \sum_j a_j$. It is implicit wherever the
map from input to features is simple: in vision, $P$ is the identity over pixels, and in language models it maps
characters to the tokens that contain them. With the rescale rule of DeepLIFT and G-DeepSHAP~\cite{shrikumar2017learning,
chen2022explaining}, it shares the idea of redistributing attribution in proportion to contribution; what is new is
obtaining $P$ for an extractor that is a program rather than a model (Section~\ref{sec:XML:Constructing}). We identify
an explanation $b \in \mathbb{R}^{|z|}$ produced directly in the problem space with $(b, 0)$, i.e., it credits nothing to
the null element.

Let $D$ be the diagonal matrix of the row sums of $P$, $D_{i,i} = \sum_j P_{i,j}$, the total share that element
$\bar{z}_i$ contributes to all features, and $D^\dagger$ its pseudoinverse ($D^\dagger_{i,i} = 1 / D_{i,i}$, or $0$
when $D_{i,i} = 0$).

\definition[\textit{Pushforward}]
\label{def:Pushforward}
{
The pushforward of a problem-space explanation $b \in \mathcal{B}$ is the feature-space explanation
$\psi(b) = P^\top D^\dagger b = a \in \mathbb{R}^{|x|}$.
}

The pushforward splits the attribution of each element among the features it produces, in proportion to
$P_{i,j} / D_{i,i}$. Reading column $j$ of $P$ as a distribution over the elements, this is Bayes' rule under a uniform
prior over the features: the pushforward is the Bayesian inverse of the pullback. Normalizing by $D^\dagger$ is what
makes it conserve attribution, $\sum_j a_j = \sum_i b_i$, whenever $b$ credits only elements that some feature reads;
crediting each feature with the plain sum of its elements' attribution instead would count a byte once for every
feature that reads it, which Section~\ref{sec:XML:PushforwardFaithfulness} shows is markedly less faithful. To our
knowledge, the pushforward is new. Summing an explanation over the regions that contain its elements, such as
superpixels or the parsed regions of a PE file~\cite{coull2019activation, bose2020explaining}, is the special case in
which every element feeds a single region; real features share their elements, and the pushforward splits each
element's attribution among all the features it produces. It is also the harder direction: it moves an explanation
backwards through a non-invertible program, into features that the explained model never computes, and so gives a
raw-byte detector a vocabulary in which its reliance recurs across files.

\begin{figure*}[htb]
\centering

% Styles
\tikzset{
    >=Stealth,
    font=\small,
    panelbox/.style={
        rounded corners,
        draw=black!60,
        very thick,
        inner sep=4mm,
        fill=white,
    },
    stage/.style={
        draw=black,
        thick,
        rounded corners,
        minimum width=3.2cm,
        minimum height=0.9cm,
        align=center,
        fill=black!3
    },
    mapbox/.style={
        draw=black,
        thick,
        rounded corners,
        minimum width=3.2cm,
        minimum height=0.9cm,
        align=center,
        fill=black!2
    },
}

\begin{tikzpicture}[
    node distance=12mm and 10mm
]

% --- Core nodes -------------------------------------------------------

% feature-space explanation a (left, top)
\node[stage, fill=black!2, minimum width=3.4cm] (anode) {
    concrete explanation\\[-1pt]
    $\displaystyle a \in \mathcal{A}$
};

% feature-space explanation a* (left, bottom; pushforward result)
\node[stage, fill=black!2, minimum width=3.4cm, below=5mm of anode] (astarnode) {
    tranferred explanation\\[-1pt]
    $\displaystyle a^{*} \in \mathcal{A}$
};

% pullback tau (middle, upper)
\node[mapbox, right=12mm of anode] (taunode) {
    pullback\\[-1pt]
    $\displaystyle \tau: \mathcal{A} \mapsto \mathcal{B}$
};

% pushforward psi (middle, lower)
\node[mapbox, right=12mm of astarnode] (psinode) {
    pushforward\\[-1pt]
    $\displaystyle \psi: \mathcal{B} \mapsto \mathcal{A}$
};

% problem-space explanation b* (right, top; pullback result)
\node[stage, fill=black!2, minimum width=3.4cm, right=12mm of taunode] (bstarnode) {
    transferred explanation\\[-1pt]
    $\displaystyle b^{*} \in \mathcal{B}$
};

% problem-space explanation b (right, bottom)
\node[stage, fill=black!2, minimum width=3.4cm, below=5mm of bstarnode] (bnode) {
    concrete explanation\\[-1pt]
    $\displaystyle b \in \mathcal{B}$
};

% Indicate closeness / correspondence within each space
\draw[<->, dashed, thick]
    (anode.south) -- node[right, font=\scriptsize] {$\approx$} (astarnode.north);

\draw[<->, dashed, thick]
    (bstarnode.south) -- node[right, font=\scriptsize] {$\approx$} (bnode.north);


% --- Arrows -----------------------------------------------------------

% Top row: a --tau--> b*
\draw[->, thick]
    (anode.east) -- (taunode.west);

\draw[->, thick]
    (taunode.east) -- (bstarnode.west);

% Bottom row: b --psi--> a*
\draw[->, thick]
    (bnode.west) -- (psinode.east);

\draw[->, thick]
    (psinode.west) -- (astarnode.east);

% Extra headroom at the top of each panel (no change to sides/bottom)
\coordinate (leftpad)  at ($(anode.north)+(0,7mm)$);
\coordinate (rightpad) at ($(bstarnode.north)+(0,7mm)$);

% --- Background region boxes (left/right only) -----------------------

\begin{scope}[on background layer]

\node[panelbox,
      fit=(anode) (astarnode) (leftpad),
      inner sep=2.5mm
] (leftpanel) {};

\node[panelbox,
      fit=(bstarnode) (bnode) (rightpad),
      inner sep=2.5mm
] (rightpanel) {};

\end{scope}

\node[font=\footnotesize,
      anchor=north,
      align=center]
    at ($(leftpanel.north)+(0,-1mm)$)
    {feature-space explanations\\$\mathcal{A} \subseteq \mathbb{R}^{|x|}$, e.g., $\mathcal{A} \subseteq \mathbb{R}^{2^{10}}$};

\node[font=\footnotesize,
      anchor=north,
      align=center]
    at ($(rightpanel.north)+(0,-1mm)$)
    {problem-space explanations\\$\mathcal{B} \subseteq \mathbb{R}^{|z|}$, e.g., $\mathcal{B} \subseteq \mathbb{R}^{2^{30}}$};

\end{tikzpicture}

\caption{Explanation Transformation. A feature-space explanation $a \in \mathcal{A}$ is pulled back via $\tau$ to yield a problem-space explanation $b^{*} \in \mathcal{B}$ close to $b$. Conversely, a problem-space explanation $b \in \mathcal{B}$ is pushed forward via $\psi$ to yield a feature-space explanation $a^{*} \in \mathcal{A}$ close to $a$.}
\label{fig:pullback-pushforward}
\end{figure*}


\subsection{Constructing the Provenance Matrix}
\label{sec:XML:Constructing}

Definition~\ref{def:ExplanationProvenance} says what $P$ must satisfy, not how to obtain it. We use two constructions,
which trade exactness and resolution against access and cost.

\paragraph{Instrumentation (white-box).} Given the extractor's source code, each feature can record the elements it is
computed from while it is computed. We adopt three conventions. An element contributes to a feature if the feature's
value is computed from it (read, counted, or hashed); indirect influence, such as a denominator that normalizes or a
pointer that says where to read, does not count. A feature that counts items gives each item one unit, spread evenly
over the item's elements, and a feature computed from content gives each element one unit. A hashed feature receives
every item that falls into its bucket. Normalizing each feature's contributions gives $P$. Instrumentation is a
one-time engineering cost per extractor, adds little to the cost of extraction, and gives $P$ at the resolution of
single elements. It is exact for features that are additive in their elements, such as counts and fields read
directly, and a convention for the others, such as hashes and entropies, which do not decompose over their elements.

\paragraph{Occlusion (black-box).} Without the source, $P$ can be estimated from the extractor's outputs alone.
Partition the elements into $K$ windows, replace each window in turn with a fill value, and re-extract the features:
the provenance of feature $j$ on window $k$ is proportional to the change $|x_j - x_j^{(k)}|$ that occluding the
window causes, spread evenly over the window's elements, and a feature that no occlusion changes goes to the null
element. The estimate costs $K + 1$ extractions per file and has the resolution of the windows. For additive features
it recovers the instrumented $P$, and for the others it measures their response to removal instead of following a
convention; but it cannot see an element whose replacement leaves the features unchanged, such as a field that
already holds the fill value, and it can create features the file does not have, such as a new low-entropy window.
